\documentclass[11pt]{article}

\usepackage[T1]{fontenc}
\usepackage[utf8]{inputenc}
\usepackage{lmodern}
\usepackage{amsmath,amssymb,amsfonts,amsthm}
\usepackage[margin=1in]{geometry}
\usepackage{microtype}
\usepackage{xcolor}
\usepackage[colorlinks=true,linkcolor=teal,urlcolor=teal]{hyperref}

\newcommand{\F}{\mathcal F_{\mathrm{orb}}}
\newcommand{\T}{\mathcal T_1}
\newcommand{\K}{\mathcal K}
\newcommand{\Bop}{\mathcal B}
\newcommand{\C}{\mathbb C}
\newcommand{\Tr}{\operatorname{Tr}}
\newcommand{\dd}{\,\mathrm d}
\newcommand{\B}{\mathrm B}

\newtheorem*{lemmaD}{Lemma D.2 (rephrased)}

\title{A Short Guide to Lemma D.2\\
\large F\o lner averaging and exact covariance in the Gaussian model}
\author{}
\date{}

\begin{document}
\maketitle

\begin{abstract}
Lemma D.2 converts an optimization averaged over larger and larger boxes of
displacements into an optimization at the origin over exactly covariant
maps.  This is a noncompact-group version of symmetrizing a channel.  The
main technical point is that averaged channels need not have a
trace-preserving limit in infinite dimension; the limit can be
subnormalized because probability may escape to arbitrarily high energy.
This note supplies the background and expands the compactness argument that
is only referenced in the paper's short proof.
\end{abstract}

\section{Why the reduction is useful}

The Gaussian cloning problem asks for a channel that performs well on every
displaced input
\[
  \Phi_z=D(z)\Phi_0D(z)^*,\qquad z\in\C^{r_d},
\]
against the correspondingly displaced target
\[
  \Psi_z=D(\sqrt\gamma z)\Psi_0D(\sqrt\gamma z)^*.
\]
The flat-prior criterion averages the root fidelity over a large cube of
$z$'s.  A general channel may behave differently at different $z$.  Lemma
D.2 shows that, in the large-cube limit, it is enough for an upper bound to
consider maps satisfying the exact covariance relation
\[
  \Gamma(D(z)XD(z)^*)
  =D(\sqrt\gamma z)\Gamma(X)D(\sqrt\gamma z)^*.
\]
For such a map, performance is the same at every displacement, so it can be
evaluated at $z=0$.

\section{Background}

\subsection{Root fidelity and symmetry averaging}

For positive trace-class operators $R,S$, possibly subnormalized, the root
fidelity is
\[
  \B(R,S)=\Tr\sqrt{R^{1/2}SR^{1/2}}.
\]
The two properties used here are
\begin{align*}
  \B(URU^*,USU^*)&=\B(R,S),\\
  \B\left(\int R_z\dd\mu(z),\int S_z\dd\mu(z)\right)
  &\geq\int\B(R_z,S_z)\dd\mu(z),
\end{align*}
for a unitary $U$ and a probability measure $\mu$.  The second line is joint
concavity.  It says that averaging conjugated copies of a channel cannot
decrease its average fidelity against the correspondingly conjugated
targets.

\subsection{F\o lner cubes}

Identify $\C^{r_d}$ with the additive group $\mathbb R^{2r_d}$, and let
\[
  \mathsf Z_L=
  \{z:|\Re z_\ell|\leq L,\ |\Im z_\ell|\leq L
  \text{ for every }\ell\}.
\]
For every fixed $v\in\C^{r_d}$,
\begin{equation}\label{eq:folner-property}
  \frac{|(\mathsf Z_L+v)\mathbin\triangle\mathsf Z_L|}
       {|\mathsf Z_L|}\longrightarrow0.
\end{equation}
Here $\triangle$ denotes symmetric difference.  This is the F\o lner
property: shifting a large cube changes only a boundary layer, whose relative
volume tends to zero.

For a compact group one could average over the whole group using normalized
Haar measure and obtain an exactly covariant channel immediately.  Phase
space is noncompact and has no invariant probability measure.  F\o lner
averaging replaces that unavailable global average by averages over expanding
cubes, followed by a compactness limit.

\subsection{Why weak-* compactness is needed}

The multimode Fock space $\F$ is infinite-dimensional.  Its trace class is
the dual of the compact operators:
\[
  \T(\F)=\K(\F)^*.
\]
Closed norm-bounded sets of trace-class operators are compact in the weak-*
topology coming from tests against compact operators, though not in trace
norm.  In this topology trace can be lost.  For instance,
\[
  |n\rangle\langle n|\stackrel{w^*}{\longrightarrow}0
  \quad\text{on }\K(\F),
\]
although every term has trace one.  This describes mass escaping to
arbitrarily high photon number.

To remember that the approximating outputs are normalized, one works first
with the unitization
\[
  \widetilde{\K}=\K(\F)\oplus\C I.
\]
A weak-* limit is a state on $\widetilde{\K}$.  Its restriction to
$\K(\F)$ is represented by a positive trace-class operator of trace at most
one; the missing trace is the component at the quotient character, or
``infinity.''  A normalized trace-class target has no component there, so
the escaped mass has zero fidelity with it.

\section{Statement of Lemma D.2}

Define the input and output actions
\begin{align}
  \alpha_z^{\mathrm{in}}(X)&=D(z)XD(z)^*,\label{eq:input-action}\\
  \alpha_z^{\mathrm{out}}(X)&=D(\sqrt\gamma z)XD(\sqrt\gamma z)^*.
  \label{eq:output-action}
\end{align}
Then $\Phi_z=\alpha_z^{\mathrm{in}}(\Phi_0)$ and
$\Psi_z=\alpha_z^{\mathrm{out}}(\Psi_0)$.

\begin{lemmaD}
For the centered cubes $\mathsf Z_L$,
\begin{equation}\label{eq:main-bound}
  \limsup_{L\to\infty}\sup_\Lambda
  \frac1{|\mathsf Z_L|}\int_{\mathsf Z_L}
  \B\bigl(\Lambda(\Phi_z),\Psi_z\bigr)\dd z
  \leq
  \sup_\Gamma\B\bigl(\Gamma(\Phi_0),\Psi_0\bigr).
\end{equation}
The left supremum is over quantum channels.  The right supremum is over
completely positive trace-nonincreasing maps satisfying
\[
  \Gamma\circ\alpha_z^{\mathrm{in}}
  =\alpha_z^{\mathrm{out}}\circ\Gamma
  \qquad(z\in\C^{r_d}).
\]
The limiting map $\Gamma$ may have subnormalized output.
\end{lemmaD}

The inequality is in the direction needed for a converse: every sequence of
nearly optimal channels on expanding cubes produces an exactly covariant
competitor on the right.

\section{Proof walkthrough}

\subsection{Step 1: average a channel over a finite cube}

For a channel $\Lambda$, define
\begin{equation}\label{eq:averaged-channel}
  \overline\Lambda_L
  =\frac1{|\mathsf Z_L|}\int_{\mathsf Z_L}
  \alpha_{-z}^{\mathrm{out}}\circ\Lambda
  \circ\alpha_z^{\mathrm{in}}\,\dd z.
\end{equation}
Every integrand is a channel, so $\overline\Lambda_L$ is a channel.  Joint
concavity and unitary invariance of root fidelity give
\begin{align}
  \B\bigl(\overline\Lambda_L(\Phi_0),\Psi_0\bigr)
  &\geq\frac1{|\mathsf Z_L|}\int_{\mathsf Z_L}
  \B\bigl(
    \alpha_{-z}^{\mathrm{out}}\Lambda(\Phi_z),\Psi_0
  \bigr)\dd z\notag\\
  &=\frac1{|\mathsf Z_L|}\int_{\mathsf Z_L}
  \B\bigl(\Lambda(\Phi_z),\Psi_z\bigr)\dd z.
  \label{eq:concavity-step}
\end{align}
Thus the original average payoff is bounded by the origin payoff of the
averaged channel.

\subsection{Step 2: extract a weak-* limit of averaged channels}

Choose $L_j\to\infty$ and channels $\Lambda_j$ whose payoffs approach the
left side of \eqref{eq:main-bound}.  Let $\overline\Lambda_j$ denote their
averages.  The adjoints restricted to the unitized compact operators are
unital completely positive contractions
\[
  \overline\Lambda_j^*:
  \widetilde{\K}\longrightarrow\Bop(\F)=\T(\F)^*.
\]
The algebra $\widetilde{\K}$ is separable.  On bounded subsets of
$\Bop(\F)$, the weak-* topology is metrizable because the predual
$\T(\F)$ is separable.  A diagonal argument on a countable dense operator
system therefore produces a subsequence and a pointwise weak-* limit
\[
  \widehat\Gamma^*:
  \widetilde{\K}\longrightarrow\Bop(\F).
\]
Positivity at every matrix level is weak-* closed, so
$\widehat\Gamma^*$ is completely positive; convergence at the identity
shows that it is unital.

Restrict $\widehat\Gamma^*$ to $\K(\F)$ and define
$\Gamma:\T(\F)\to\T(\F)$ by
\begin{equation}\label{eq:limit-duality}
  \Tr[\Gamma(X)A]=\Tr[X\widehat\Gamma^*(A)],
  \qquad A\in\K(\F).
\end{equation}
This gives a completely positive map.  Let $(e_n)$ be increasing
finite-rank projections converging strongly to $I$.  For $X\geq0$,
\begin{align*}
  \Tr\Gamma(X)
  &=\sup_n\Tr[\Gamma(X)e_n]\\
  &=\sup_n\Tr[X\widehat\Gamma^*(e_n)]
  \leq\Tr[X\widehat\Gamma^*(I)]
  =\Tr X.
\end{align*}
Hence $\Gamma$ is trace nonincreasing.  Equality need not hold because the
restriction from $\widetilde{\K}$ to $\K(\F)$ discards possible mass at
infinity.

\subsection{Step 3: the F\o lner property gives exact covariance}

For fixed $v\in\C^{r_d}$, changing variables in
\eqref{eq:averaged-channel} gives
\begin{equation}\label{eq:boundary-estimate}
  \left\|
  \overline\Lambda_L(\alpha_v^{\mathrm{in}}X)
  -\alpha_v^{\mathrm{out}}\overline\Lambda_L(X)
  \right\|_1
  \leq
  \frac{|(\mathsf Z_L+v)\mathbin\triangle\mathsf Z_L|}
       {|\mathsf Z_L|}\,\|X\|_1.
\end{equation}
The integrands are trace-norm contractions, and only the two boundary pieces
left by shifting the cube contribute.  By \eqref{eq:folner-property}, the
right side tends to zero.  Pairing \eqref{eq:boundary-estimate} with compact
operators and passing to the extracted weak-* limit yields
\[
  \Gamma(\alpha_v^{\mathrm{in}}X)
  =\alpha_v^{\mathrm{out}}\Gamma(X).
\]
This holds for every fixed $v$.  One may first prove it on a countable dense
set of $v$'s used in the diagonal extraction; strong continuity of the Weyl
operators, which gives trace-norm continuity of their action on trace class,
then extends it to all $v$.

\subsection{Step 4: pass the fidelity payoff to the limit}

Put $\rho_j=\overline\Lambda_j(\Phi_0)$.  Viewed as states on
$\widetilde{\K}$, these converge weak-* to a state
$\widehat\rho$.  Its restriction to $\K(\F)$ is the trace-class operator
\[
  \rho=\Gamma(\Phi_0),
\]
which may have trace less than one.

For positive functionals on a unital $C^*$-algebra, the
Alberti--Uhlmann variational formula is
\begin{equation}\label{eq:AU}
  \B(\phi,\psi)
  =\inf_{A>0\ \mathrm{invertible}}
  \frac12\bigl(\phi(A)+\psi(A^{-1})\bigr).
\end{equation}
For fixed $\psi$, the right side is an infimum of weak-* continuous affine
functions of $\phi$, so it is weak-* upper semicontinuous.  Therefore
\[
  \limsup_j\B(\rho_j,\Psi_0)
  \leq\B_{\widetilde{\K}}(\widehat\rho,\widehat\Psi_0).
\]
The target $\Psi_0$ is normalized and supported on the original compact
operator algebra.  Its extension has zero component at infinity.  The
central-summand, or mass-at-infinity, identity therefore gives
\[
  \B_{\widetilde{\K}}(\widehat\rho,\widehat\Psi_0)
  =\B\bigl(\Gamma(\Phi_0),\Psi_0\bigr).
\]
Combining this inequality with \eqref{eq:concavity-step}, and taking the
supremum over all possible covariant limits $\Gamma$, proves
\eqref{eq:main-bound}.

\section{What subnormalization means}

Subnormalization is not an extra physical resource and does not enlarge the
original class of channels before taking the limit.  Every
$\overline\Lambda_j$ is trace preserving.  The trace defect appears only
because compact observables cannot detect output weight that moves to
unbounded photon number.  Such weight has vanishing overlap with the fixed
trace-class target, so retaining it would not improve the limiting fidelity.

The elementary sequence $|j\rangle\langle j|$ illustrates the point.  Its
weak-* limit on compact operators is zero, while on the unitization the
missing unit trace survives entirely at the quotient character.  Against a
fixed normalized target in the original Fock space, that quotient component
is orthogonal and contributes no fidelity.

\section{Dependencies and role in the paper}

Lemma D.2 is the purely quantum specialization of Lemma C.2.  Relative to
that hybrid result, the classical coordinate is omitted and the translation
group is only $\C^{r_d}$.  Its proof uses:
\begin{enumerate}
  \item joint concavity and unitary invariance of root fidelity;
  \item the F\o lner property of Euclidean cubes;
  \item weak-* compactness for completely positive maps on the unitized
  compact-operator algebra;
  \item upper semicontinuity of transition amplitude and the
  mass-at-infinity identity explained in Lemma C.1.
\end{enumerate}
It does not use Lemma D.1.  Lemma D.1 enters the separate stochastic-order
chain beginning with Lemma D.3.

In the proof of Theorem 3.4, the product quantum-limited amplifier supplies
the lower bound.  For the converse, Lemma D.2 reduces the large flat-prior
average to
\[
  \B\bigl(\Gamma(\Phi_0),\Psi_0\bigr)
\]
for an exactly displacement-covariant trace-nonincreasing map.  Phase
twirling then makes the map phase covariant without lowering this origin
fidelity, and Lemma D.5 supplies the explicit multimode upper bound.  Thus
D.2 is the symmetry-and-compactness step that makes the later photon-number
argument applicable.

\section{Minimal prerequisite checklist}

The conceptual core requires only group averaging, concavity of fidelity,
and the fact that large cubes have negligible relative boundary.  The
functional-analytic part additionally uses trace-class/compact-operator
duality, weak-* compactness, unitization, and upper semicontinuity of the
Alberti--Uhlmann transition amplitude.  Its practical message is simple:
an asymptotically translation-invariant average may be replaced, for an
upper bound, by one exactly covariant map evaluated at the origin.

\end{document}
